\begin{figure}[ht]
\centering
\begin{tikzpicture}[scale=.83,>=Latex]
 \draw[thick] (0,0) circle (1.05);
 \node at (0,0) {$\D$};
 \fill (1.05,0) circle (2pt);
 \node[below right] at (1.05,0) {$\lambda$};
 \draw[->] (1.05,0) -- (1.8,0);
 \draw[->] (2.1,.5) to[bend left=15] node[above] {$h$} (4,.5);
 \draw[thick] (5.35,0) ellipse (1.1 and .85);
 \node at (5.25,-.15) {$\Omega$};
 \draw[dashed] (6.45,-1.3) -- (6.45,1.3);
 \fill (6.45,0) circle (2pt);
 \draw[->,thick] (6.45,0) -- (7.6,0) node[right] {$q(\lambda)$};
 \node[above right] at (6.45,.08) {$h(\lambda)$};
 \fill (4.95,.35) circle (1.5pt);
 \node[above] at (4.95,.35) {$z$};
\end{tikzpicture}
\caption{The map $h$ sends the exterior of the circle to the exterior of the domain; the supporting line is shown schematically. Increasing the circle radius gives the outward normal $q(\lambda)$; convexity places every $z\in\overline\Omega$ on the inner side of the line.}
\label{fig:normal}
\end{figure}
